\par\SourceLeaf{434}\Needspace{7\baselineskip}\bigskip
\begin{center}
\pdfbookmark[2]{Gauss an Olbers, Göttingen, 24. Julius 1816}{pallas-brief-18160724}
\textsc{Gauss} an \textsc{Olbers}, Göttingen, 24. Julius 1816.
\end{center}
\noindent\hspace*{\parindent}. . . . .\quad Die heurige Pallas-Opposition [XI.] verträgt sich noch immer sehr gut mit meiner Theorie und bestätigt die Nothwendigkeit einer Vergrösserung der Jupitersmasse. Ich finde jetzt \(\frac{\text{\AstroSymbols ☉}}{\text{\AstroSymbols ♃}}=1050\) und die wahrscheinliche Ungewissheit dieser Zahl sehr nahe \(=1\); \textsc{Laplace} findet aus der Saturnsbewegung 1070 mit einer etwas grössern Ungewissheit; dieser Unterschied ist also ganz enorm; allein ich glaube nicht, dass man daraus berechtigt ist, auf eine verschiedene affinitas chemica zu schliessen, da \textsc{Laplaces} Resultat sich auf \textls[180]{seine} Saturnstheorie gründet, die nach einer Methode entwickelt ist, deren Zulänglichkeit bezweifelt werden kann. \textls[180]{Meine} Bestimmung der Jupitersmasse wünsche ich übrigens noch nicht bekannt zu sehen. Jede neu hinzugekommene Pallasopposition wird die Genauigkeit vergrössern. Freilich ungeheure Arbeit, da die Rechnung in 14 Tagen für Eine Opposition sich nicht machen lässt. Sehr neugierig bin ich darauf, was für Resultate Juno demnächst geben wird.\quad . . . . .
\par
